\documentclass[11pt]{amsart}
\usepackage[T1]{fontenc}
\usepackage[utf8]{inputenc}
\usepackage{amsmath,amssymb,amsthm,mathtools}
\usepackage{microtype}
\usepackage[colorlinks=true,linkcolor=blue,citecolor=blue,urlcolor=blue]{hyperref}

\newtheorem{theorem}{Theorem}
\newtheorem{lemma}[theorem]{Lemma}

\theoremstyle{remark}
\newtheorem{remark}[theorem]{Remark}

\newcommand{\Cyc}[1]{C_{#1}}
\newcommand{\Z}{\mathbb{Z}}
\newcommand{\ord}{\operatorname{ord}}


\begin{document}

\title[No Abelian $(64681,441,3)$-Difference Set]{No Abelian $(64681,441,3)$-Difference Set Exists\thanks{AutoMath assisted in generating the solution and preparing this draft.}}
\date{April 15, 2026}

\begin{abstract}
We show that no abelian group of order $64681$ admits a $(64681,441,3)$-difference set.
The proof is short: since $64681=71\cdot 911$ is squarefree, every abelian group of that order is cyclic.
The first multiplier theorem makes $73$ a multiplier, and after translation any putative difference set is invariant under multiplication by $73$ on $\Cyc{64681}$.
Every nonzero $73$-orbit has size $35$, so an invariant set has size $\delta+35a$ with $\delta\in\{0,1\}$.
This cannot equal $441$.
\end{abstract}

\maketitle

\section{Introduction}
Gordon's 2022 analysis of abelian difference sets with small $\lambda$ leaves $(64681,441,3)$ on the residual list of open abelian $\lambda=3$ cases \cite{Gordon2022}.
We record a direct multiplier-orbit obstruction that removes this tuple.

\begin{theorem}
No abelian group of order $64681$ admits a $(64681,441,3)$-difference set.
\end{theorem}

\begin{proof}
Let $G$ be an abelian group of order
\[
64681 = 71\cdot 911.
\]
Since the order is squarefree, $G$ is cyclic, so it is enough to rule out a difference set in $\Cyc{64681}$.

Assume for contradiction that $D\subseteq \Cyc{64681}$ is a $(64681,441,3)$-difference set.
Its order is
\[
n=k-\lambda=441-3=438=2\cdot 3\cdot 73.
\]
Because $73$ is a prime divisor of $n$ with $73>\lambda$ and $\gcd(73,64681)=1$, the first multiplier theorem implies that $73$ is a numerical multiplier of $D$.
Thus there exists $s\in \Cyc{64681}$ such that
\[
73D=D+s.
\]
Since
\[
\gcd(73-1,64681)=\gcd(72,64681)=1,
\]
we may choose $g\in \Cyc{64681}$ with $(73-1)g=-s$.
Then the translate $D'=D+g$ satisfies
\[
73D'=73D+73g=D+s+73g=D+g=D'.
\]
Replacing $D$ by $D'$, we may assume from the start that $73D=D$.

We now compute the orbit sizes of the map $x\mapsto 73x$ on $\Cyc{64681}$.
The point $0$ is fixed.
If $x$ is a nonzero multiple of $911$, then its orbit size is $\ord_{71}(73)$.
Since $73\equiv 2\pmod{71}$ and $2^{35}\equiv 1\pmod{71}$, we have $\ord_{71}(73)=35$.
Likewise, if $x$ is a nonzero multiple of $71$, then the orbit size is $\ord_{911}(73)$, and a direct computation gives $\ord_{911}(73)=35$.
For a unit $x$, the orbit size is
\[
\operatorname{lcm}\bigl(\ord_{71}(73),\ord_{911}(73)\bigr)=35.
\]
Therefore every nonzero orbit has size $35$.

Because $D$ is $73$-invariant, it is a union of $73$-orbits.
Hence
\[
|D|=\delta+35a
\]
for some integer $a\ge 0$ and some $\delta\in\{0,1\}$ according to whether $0\in D$.
But $|D|=441$, and neither $441$ nor $440$ is divisible by $35$.
This contradiction proves that no such difference set exists.
\end{proof}

\begin{remark}
The obstruction is entirely orbit-theoretic: once the valid multiplier $73$ is chosen, the orbit-size semigroup on $\Cyc{64681}$ misses the required weight $441$.
\end{remark}

\begin{thebibliography}{99}

\bibitem{Gordon2022}
Daniel M.~Gordon,
\emph{On difference sets with small lambda},
J. Algebraic Combin. \textbf{55} (2022).

\bibitem{Lander}
E.~S.~Lander,
\emph{Symmetric Designs: An Algebraic Approach},
Cambridge University Press, 1983.

\end{thebibliography}

\end{document}
